% This is the chapter 'Thermodynamics' for the Hodgkin-Huxley neuron model V2.0 book
% Last edited 2026.06.22

\ifx\magyar\undefined
\chapter{Thermodynamics for neurons\label{ch:Thermodynamics}}
\else
\chapter[A neuronok termodinamikája]{A neuronok termodinamikája\label{ch:Thermodynamics}}
\fi
\iflatexml
\else
\minitoc
\fi

A termodinamika szokásos értelmezése izotróp és homogén szabad térben elhelyezkedő, nagy számú olyan részecskére vonatkozik, amelyek csak közvetlen ütközéssel tudnak kölcsönhatásba lépni. Mint másutt részleteztük, a neuronok esetében általában egyik feltétel sem teljesül, ezért a termodinamikai következtetéseket újra kell származtatni és azok érvényességét megvizsgálni. A termodinamika alapvető összefüggései azonban érvényben maradnak. Ebben a fejezetben a termodinamika  olyan vonatkozásait emeljük ki, amelyek általában elektrolitok esetén, de különösen neuronok esetén jelentősen eltérnek a megszokottól.

\section[Termodinamikus tér]{A "termodinamikai elektromos tér" \label{sec:Physics-ThermodynamicField}}
\fi


Szegmensekre osztott elektrolitban az elektromos töltések kiegyensúlyozottak; megfelelő fizikai körülmények között
azonban kiegyensúlyozatlanná is válhatnak.
Egy, a biológia számára fontos eset, amikor egy féligáteresztő hártya
választ el két, nagyon eltérő koncentrációjú ionos oldatokat
tartalmazó szegmenst. Az eltérő koncentrációk 
a jól ismert Nernst-Planck összefüggésnek megfelelően feszültséget generálnak, amint azt a 
\hyperlink{phys:Nernst-Planck}
{\textbf{Physics panel~\ref{phys:Nernst-Planck}}} tárgyalja.
A nyugalmi potenciál esetére vonatkozó számítási módot a \hyperlink{phys:EquivalentField_Membrane}
{\textbf{Physics panel~\ref{phys:EquivalentField_Membrane}}},
annak eredményét pedig a 
\hyperlink{num:EquivalentField_Membrane}
{\textbf{Numerics panel~\ref{num:EquivalentField_Membrane}}} mutatja.

Egyszerűen fogalmazva, a Nernst-Planck egyenlet azt állítja,
hogy az ionok koncentrációjának változása változást okoz a
elektromos térerőben (és megfordítva);
továbbá, hogy stacionárius esetben azok változatlanok maradnak.
Vegyük észre, hogy ionkeverék esetében közös, töltés-generált,
elektromos erőtér van, és az minden ionra azonos, míg a közös
elektromos erőtér az egyes ionoknál eltérő koncentrációkat generál. 
Megjegyezzük, hogy a neuronon belül a gradiensek élesen változnak,
így ott csupán a differenciális alak alkalmazható;
a statikus makroszkopikus integrális alak használata félrevezető.

\physicspanel{phys:Nernst-Planck}{A Nernst-Planck törvény}
{
The Nernst-Planck electrodiffusion equation in a balanced state
\begin{equation}
	\frac{d}{dz}V_{m}(z)=-\frac{RT}{q*F}\frac{1}{C_{k}(z)}\frac{d}{dz}C_{k}(z)\label{eq:NernstPlanck}
\end{equation}
describes in one dimension the mutual dependence of the \emph{spatial gradients} of the
electrical and thermodynamic fields on each other. In good textbooks (see, for example,~\cite{KochBiophysics:1999},
Eq (11.28)), its derivation is exhaustively detailed. In the
equation, $z$ is the spatial variable across the direction of the
changed invasion parameter, $R$ is the gas constant, $F$ is the
Faraday's constant, $T$ is the temperature, $q$ the valence
of the ion, ${V_m(z)}$ the potential, and ${C_k(z)}$ the concentration
of the chemical ion. The two sides are the derivatives of the equation
\begin{equation}
	V_{m} =\frac{RT}{q*F}\ln{\biggl(\frac{C_{k}^{ext}}{C_{k}^{int}}\biggr)}\label{eq:Nernst1}
\end{equation}
\noindent known as \textit{Nernst equation}.  In other words, \textbf{Eqs.(\ref{eq:NernstPlanck}) and Eq.(\ref{eq:Nernst1}) are the differential and integral formulations of the same knowledge}. The limits of the integration are chosen arbitrarily (by choosing the reference concentration $C_{k}^{int}$).
}





